\documentclass[11pt]{article}
\usepackage[margin=1in]{geometry}
\usepackage{amsmath, amssymb, amsthm}
\usepackage{booktabs}
\usepackage{array}
\usepackage{siunitx}
\usepackage{hyperref}
\usepackage{xcolor}
\usepackage{listings}

\hypersetup{colorlinks=true, linkcolor=blue, urlcolor=blue}

\title{BCM $K_{\text{rut}}$ Sweep --- Observations from the 10-Gene Test Panel}
\author{Net-Seq-Modeling project}
\date{\today}

\newcommand{\Krut}{K_{\text{rut}}}
\newcommand{\kload}{k_{\text{load}}}
\newcommand{\kribo}{k_{\text{ribo}}}
\newcommand{\Dstar}{D^{\!*}}
\newcommand{\DstarWT}{\Dstar_{\text{WT}}}
\newcommand{\Snorm}{S_{\text{norm}}}
\newcommand{\Sexp}{S_{\text{exp}}}
\newcommand{\SWT}{S_{\text{WT}}}
\newcommand{\SBCM}{S_{\text{BCM}}}

\begin{document}
\maketitle

\begin{abstract}
We report the results of a 10-gene test panel for a $\Krut$ (Rho loading rate)
sweep against bicyclomycin (BCM) NET-seq data, using intrinsic dwell-time
profiles $\Dstar_{\text{WT}}(x)$ obtained from CMA-ES deconvolution of wildtype
NET-seq. The sweep exercises a per-candidate $\Krut$ refactor of the GPU TASEP
kernel that enables a 50-point grid evaluation in a single batched launch. We
observe that $\Dstar_{\text{WT}}$ does not cleanly extrapolate to the BCM
condition for most panel genes under the current model: nearly all 10 genes
find a best-fit $K^{*}_{\text{rut}}$ well below the WT-assumed value of $0.13$. Two genes
pile at the $\Krut=0$ grid boundary (\texttt{rpoB}, \texttt{yihO}), and
\texttt{yihO} exhibits a residual MSE an order of magnitude larger than any
other gene, flagging it as a model-failure candidate for future joint re-fit.
A $\Delta$-signal metric, designed to cancel $\Dstar_{\text{WT}}$ errors shared
between conditions, reports ``no rho change predicted'' for 7/10 panel genes at
$n_{\text{runs}}=5000$; two genes (\texttt{aceA}, \texttt{rpoB}) now reveal
interior $\Delta$-metric minima at $K_{\text{rut}} \approx 0.06$.
\end{abstract}

\section{Setup}

The TASEP forward model produces a per-position NET-seq-like density
$S(x)$ given an intrinsic dwell-time profile $D(x)$ and a set of kinetic
parameters $\theta = (\kload, \kribo, \Krut, w_{\text{RNAP}}, \ldots)$:
\begin{equation}
  S(x) \;=\; \mathcal{T}\bigl[D;\;\theta\bigr](x).
\end{equation}
Here $\mathcal{T}$ is the nonlinear TASEP operator that emerges from
first-passage dynamics with hard-body exclusion, Rho-mediated premature
termination, and ribosome coupling. Although one might be tempted to write
$S = D \ast j$ for some flux kernel $j$, the underlying dynamics are not
linear in $D$ (see \S\ref{sec:math} for the structure). The correct
formulation is a nonlinear functional, pointwise of the form
$S(x) \propto D(x)\cdot j(x;D,\theta)$.

For each gene, the CMA-ES deconvolution solves
\begin{equation}
  \DstarWT \;=\; \arg\min_{D} \;\bigl\lVert\mathcal{T}[D;\;\Krut=0.13,\;\kribo]
  \;-\; \SWT\bigr\rVert^2,
\end{equation}
with $\kribo = \text{TE}_{\text{gene}}/5$ taken from
\texttt{Ecoli\_gene\_TE.csv} at load time and $\Krut$ held fixed. The
identifiability problem we address here: when $\Krut=0.13$ does not
match true in-vivo Rho activity for a given gene, the CMA-ES fit has no
free handle except $D(x)$, and any mismatch is absorbed into $\DstarWT$.

BCM is a Rho inhibitor. NET-seq under BCM is approximately equivalent to
running the same $\DstarWT$ with $\Krut \to 0$ (or a small residual value
reflecting partial inhibition). This provides the constraint that WT
alone cannot: varying $\Krut$ and comparing against $\SBCM$ identifies
whether the model can reconcile both conditions with a single $\DstarWT$.

\section{Mathematical structure of the NET-seq signal}
\label{sec:math}

\subsection{Pointwise density equals dwell times flux}

At steady state, in a mean-field approximation,
\begin{equation}
  S(x) \;\propto\; \rho(x) \;=\; D(x)\,j(x;D,\theta),
\end{equation}
i.e.\ density at position $x$ is the local dwell-time multiplied by the
flux of RNAPs that actually reach $x$.

\subsection{The flux decomposes into a loading source and two survival factors}

\begin{equation}
  j(x;D,\theta) \;=\; \kload \cdot \Pi_{\rho}(x;D,\Krut,\kribo)
                   \cdot \Pi_{\text{traffic}}(x;D,\kload)
\end{equation}

where $\Pi_{\rho}$ is the probability of surviving Rho termination from
the promoter to $x$, and $\Pi_{\text{traffic}}$ is a collision-blocking
factor from hard-body exclusion.

\paragraph{Rho survival} The Rho-loading rate at position $x$ scales
with uncovered nascent RNA length:
\begin{equation}
  \lambda_{\rho}(x;\Krut,\kribo,D) \;=\; \frac{\Krut}{L_{\text{gene}}}
    \cdot L_{\text{uncovered}}(x;\kribo,D),
\end{equation}
where $L_{\text{uncovered}}$ is the stretch of RNA between the ribosome
position $r(x)$ and the RNAP's 5$^\prime$ footprint,
\(
L_{\text{uncovered}}(x) \approx x - w_{\text{RNAP}} - 30 - r(x).
\)
The survival factor from $0$ to $x$ is
\begin{equation}
  \Pi_{\rho}(x) \;=\; \exp\!\left(-\int_{0}^{x}
    \lambda_{\rho}(x')\,D(x')\,\mathrm{d}x'\right).
\end{equation}
Slow dwell ($D$ large) increases exposure time and hence the effective
rho-loading opportunity --- so $\Pi_{\rho}$ depends on \emph{all} of
$D$ upstream of $x$, not just locally.

\paragraph{Traffic throughput} Hard-body exclusion gives an effective
velocity
\begin{equation}
  v_{\text{eff}}(x) \;=\; \frac{1}{D(x)}\bigl(1 - a\,\rho(x+1)\bigr),
\end{equation}
so the steady-state continuity with Rho attrition is
\begin{equation}
  \frac{\partial}{\partial x}\!\left[\rho(x)\,v_{\text{eff}}(x)\right]
  \;=\; -\,\lambda_{\rho}(x)\,\rho(x).
\end{equation}
This is a self-consistent nonlinear constraint; $\rho$ depends on
$v_{\text{eff}}$, which depends on $\rho(x+1)$. No closed-form inversion
exists, which is why the TASEP simulator is run numerically rather than
replaced by an analytic convolution.

\subsection{Full functional form}

Collecting:
\begin{equation}
\boxed{
  S(x) \;=\; \alpha \cdot D(x) \cdot \kload
      \cdot \exp\!\left(-\int_{0}^{x}\frac{\Krut}{L}\,
          L_{\text{uncovered}}(x';\kribo)\,D(x')\,\mathrm{d}x'\right)
      \cdot \Pi_{\text{traffic}}(x;D,\kload)
}
\end{equation}
The dependence on $D$ is local through the $D(x)$ prefactor but
non-local through both the rho-survival integrand and the implicit
continuity equation for $\Pi_{\text{traffic}}$.

\section{Sweep design}

\begin{enumerate}
\item For each gene, load $\DstarWT$ and $\kribo$ from
  \texttt{ecoli/cmaes/\{gene\}.pkl}. Load WT $S_{\text{exp}}$ from the
  same pkl (preserving the preprocessing used during the CMA-ES fit).
\item Load BCM NET-seq from
  \texttt{NETSEQ\_gene\_nomodifications\_BCM/NETSEQ\_\{gene\}.csv} and
  apply a pseudo-count to match WT preprocessing:
  \begin{equation}
    \SBCM(x) \;=\; \frac{\text{raw}(x) + q_{\text{BCM}}}
                      {\mathrm{mean}(\text{raw} + q_{\text{BCM}})},
  \end{equation}
  with $q_{\text{BCM}}$ the smallest nonzero value in the raw profile
  (observed $\approx 0.293086$ for every panel gene, consistent with a
  single global library normalization constant).
\item Evaluate 50 uniform $\Krut$ values in $[0, 0.13]$ in a single
  batched GPU kernel launch (enabled by promoting \texttt{pt\_percent}
  from a scalar to a per-candidate array
  \texttt{pt\_percent\_arr}, mirroring the existing per-candidate
  \texttt{k\_ribo\_loading\_arr} pattern).
\item Record two loss metrics for each $\Krut$, both on
  mean-normalized profiles:
  \begin{align}
    \text{MSE}_{\text{pw}}(\Krut) &=
      \frac{1}{L}\sum_{x}\bigl(\Snorm(\Krut, x) - \SBCM(x)\bigr)^2, \\
    \text{MSE}_{\Delta}(\Krut) &=
      \frac{1}{L}\sum_{x}\Bigl[\bigl(\Snorm(\Krut, x) - \Snorm(0.13, x)\bigr)
      - \bigl(\SBCM(x) - \SWT(x)\bigr)\Bigr]^2.
  \end{align}
\item Simulate at $n_{\text{runs}} = 5000$ and $\mathrm{d}t = 0.1$.
\item Save the full 50-point grid, both metric curves, and the reference
  signals per gene; compute argmin values downstream in analysis. No
  argmin is baked into the per-gene pickle.
\end{enumerate}

\section{Kernel refactor and validation}

\paragraph{GPU} \texttt{pt\_percent} is now per-candidate. Single
launch evaluates all 50 grid points at $\sim 21$ seconds per gene for
panel-average gene length. A regression test on \texttt{aceA} comparing
pre- and post-refactor flux output via \texttt{postprocess\_flux.py}
gave Pearson $r = 1.000000$ across all 1305 positions; the residual
$\approx 20\%$ magnitude offset is a pure global scale shift from
$\text{flux}[0]$ being sensitive to 1-ULP FMA reordering across kernel
recompilation, and is irrelevant for mean-normalized comparisons.

\paragraph{CPU parity} \texttt{validate\_gpu\_cpu.py} was extended to
sweep $\Krut \in \{0.0, 0.065, 0.13\}$. Pearson correlation between
CPU and GPU mean NET-seq profiles exceeds $0.998$ at all three values;
overall signal totals agree within $0.5\%$. The per-position strict
$<5\%$ tolerance fails on low-signal pixels (pre-existing Poisson-noise
regime issue), but distributional agreement is well-established.

\paragraph{CPU kernel} Not refactored by design decision. The CPU
reference \texttt{netseq\_tasep\_fast.py::\_tasep\_core} is
single-run-scalar; per-candidate $\Krut$ flows naturally via
per-invocation \texttt{parameters["KRutLoading"]} at the orchestrator
level.

\section{10-gene panel results}
\label{sec:results}

\begin{table}[h]
\centering
\small
\begin{tabular}{lllrrrrrr}
\toprule
Gene & Class & $L$ (bp) & $K^{*,\text{pw}}_{\text{rut}}$ & $K^{*,\Delta}_{\text{rut}}$ &
MSE$_{\text{pw}}$ & MSE$_{\Delta}$ & Endpoint pw & Endpoint $\Delta$ \\
\midrule
aceA & prior-bench   & 1305 & 0.032 & 0.061 & 2.03 & 1.18 & -- & -- \\
dnaK & prior-bench   & 1917 & 0.029 & 0.130 & 2.45 & 0.65 & -- & max \\
rpoB & prior-bench   & 4029 & 0.000 & 0.064 & 3.70 & 0.66 & min & -- \\
talB & prior-bench   &  954 & 0.024 & 0.130 & 2.30 & 0.76 & -- & max \\
insQ & prior-bench   & 1149 & 0.127 & 0.130 & 3.91 & 1.82 & -- & max \\
\addlinespace
rhsA & rho-dependent & 4134 & 0.005 & 0.130 & 3.48 & 1.62 & -- & max \\
rhsB & rho-dependent & 4236 & 0.027 & 0.130 & 3.14 & 0.88 & -- & max \\
ydbA & rho-dependent & 8622 & 0.035 & 0.130 & 5.12 & 2.00 & -- & max \\
yihO & rho-dependent & 1404 & 0.000 & 0.000 & 25.97 & 15.50 & min & min \\
caiT & rho-dependent & 1515 & 0.024 & 0.130 & 2.90 & 1.54 & -- & max \\
\bottomrule
\end{tabular}
\caption{Per-gene best-fit $\Krut$ under the two metrics at
$n_{\text{runs}}=5000$. Endpoint columns flag argmin landing at the
grid boundaries ($\Krut=0$ or $\Krut=0.13$), an indication that the
true minimum may lie beyond the swept range.}
\label{tab:panel}
\end{table}

\section{Observations}

\subsection{Rho-dependent genes behave as predicted}

All five known-rho-dependent genes (\texttt{rhsA}, \texttt{rhsB},
\texttt{ydbA}, \texttt{yihO}, \texttt{caiT}) have
$K^{*,\text{pw}}_{\text{rut}} \le 0.035$, substantially below the WT-assumed
$0.13$. The sweep correctly recovers the Rho-inhibition signature from
BCM using a model trained only on WT data.

\subsection{Prior-benchmark genes also shift well below $0.13$}

Only \texttt{insQ} remains near the WT value ($K^{*,\text{pw}}_{\text{rut}} =
0.127$). The other four prior-benchmark genes --- \texttt{aceA},
\texttt{dnaK}, \texttt{rpoB}, \texttt{talB} --- prefer
$K^{*,\text{pw}}_{\text{rut}} \le 0.032$. The original 5-gene benchmark was
curated for CMA-ES validation, not Rho independence; the panel is not
as clean a classifier as initially assumed. Under the current model,
these genes either (a) truly have effective in-vivo $\Krut$ below
$0.13$, or (b) have systematic $\DstarWT$ errors that the sweep
compensates for nonspecifically (e.g.\ by reducing 3$^\prime$ flux).

\subsection{Two genes pile at the $\Krut=0$ floor}

\texttt{rpoB} and \texttt{yihO} both minimize at the grid boundary. The
true minimum lies at or below $\Krut=0$, which is physically
inaccessible: the model has no handle lower than ``no Rho at all.'' The
interpretation is that $\DstarWT$ for these genes absorbed substantial
Rho-termination shape during the WT fit that the forward sweep cannot
now fully undo by reducing $\Krut$. \texttt{rhsA}, which appeared at
the floor in the $n_{\text{runs}}=500$ run, resolves to a shallow
minimum at $K^{*,\text{pw}}_{\text{rut}} = 0.005$ with higher statistics
--- still strongly Rho-dependent but no longer strictly
floor-pinned. This is exactly the identifiability artifact BCM data was
intended to diagnose.

\subsection{\texttt{yihO} is a model-failure candidate}

$\text{MSE}_{\text{pw,min}}(\texttt{yihO}) = 25.97$, an order of
magnitude larger than any other gene in the panel. No $\Krut$ in the
grid reconciles $\SWT$ and $\SBCM$ under the current model. Either the
simple kinetic Rho model is too crude to capture real termination for
this gene, there is un-modeled Rho-dependent biology (NusG antagonism,
secondary structure), or the experimental BCM condition had
gene-specific confounds. \texttt{yihO} is the primary motivation for a
joint (WT $+$ BCM) re-fit of $D$ in a follow-up change.

\subsection{The $\Delta$-metric is MC-noise-limited at
  $n_{\text{runs}}=500$}

$K^{*,\Delta}_{\text{rut}} = 0.13$ for 7/10 panel genes at
$n_{\text{runs}} = 5000$. The $\Delta$-metric computes
\(\Snorm(\Krut) - \Snorm(0.13)\), the difference of two independent
stochastic realizations, which carries $\sqrt 2$ times the Monte Carlo
noise of a single simulation. At $n_{\text{runs}} = 5000$, the
per-position noise floor in the $\Delta$-metric drops to
$\sim 1/\sqrt{5000} \cdot \sqrt{2} \approx 0.02$, roughly 3$\times$
lower than at $n_{\text{runs}} = 500$.

This reduction is sufficient to reveal interior $\Delta$-metric minima
for two genes: \texttt{aceA} ($K^{*,\Delta}_{\text{rut}} = 0.061$,
$\text{MSE}_\Delta = 1.18$) and \texttt{rpoB}
($K^{*,\Delta}_{\text{rut}} = 0.064$, $\text{MSE}_\Delta = 0.66$).
Both minima are well below the WT-assumed $0.13$, consistent with their
pointwise estimates, and confirm that the $\Delta$-signal for these
genes exceeds the noise floor at this run count.

For the remaining 7 genes the $\Delta$-metric minimum is still pinned at
$\Krut=0.13$ where $\Snorm(0.13) - \Snorm(0.13) = 0$ by construction.
The true Rho-induced shape change for these genes remains below the
$n_{\text{runs}} = 5000$ noise floor. Only \texttt{yihO} has a
$\SBCM - \SWT$ large enough to produce an unambiguous $\Delta$-metric
minimum at $\Krut = 0$.

\subsection{Summary classification}

Combining both metrics, the 10 genes partition approximately as:

\begin{itemize}
\item \textbf{$\DstarWT$ was probably correctly calibrated for
  $\Krut=0.13$}: \texttt{insQ} (pointwise near WT endpoint at 0.127;
  $\Delta$-metric also pinned at 0.13; shallow curves).
\item \textbf{$\DstarWT$ absorbed some Rho shape, $\Delta$-metric now
  informative}: \texttt{aceA}, \texttt{rpoB} (pointwise interior
  minimum well below 0.13; $\Delta$-metric reveals interior minimum at
  $\Krut \approx 0.06$ at $n_{\text{runs}}=5000$).
\item \textbf{$\DstarWT$ absorbed some Rho shape, $\Delta$-metric still
  noise-limited}: \texttt{dnaK}, \texttt{talB}, \texttt{rhsB},
  \texttt{ydbA}, \texttt{caiT} (pointwise clearly prefers $\Krut <
  0.13$ with reasonable MSE; $\Delta$-metric pinned at 0.13).
\item \textbf{$\DstarWT$ absorbed substantial Rho shape}: \texttt{rhsA}
  (pointwise minimum at $\Krut = 0.005$, near floor),
  \texttt{rpoB} (pointwise at floor).
\item \textbf{Model failure}: \texttt{yihO} (floor-pinned pointwise;
  residual MSE $> 8\times$ panel median after the sweep).
\end{itemize}

\section{Implications for the broader pipeline}

\subsection{What the sweep resolves}

For genes where both metrics agree on a $K^{*}_{\text{rut}} < 0.13$ with a clean
minimum and small residual MSE, the WT-fit $\DstarWT$ overestimates
intrinsic dwell at 3$^\prime$ regions by absorbing Rho-termination
shape. A corrected estimate would be $\DstarWT(x) \cdot$ shape
correction inferred from the BCM fit. The sweep itself does not
produce a corrected $D$; it produces a diagnostic of how far off the
WT fit is.

\subsection{What the sweep does not resolve}

The sweep cannot distinguish between (i) $\DstarWT$ absorbed Rho
shape and (ii) $\DstarWT$ absorbed some other un-modeled effect that
happens to have similar 3$^\prime$ shape. To distinguish these, a
joint re-fit is required: optimize $D$ against both $\SWT$ and
$\SBCM$ simultaneously, with $\Krut$ taking different values in each
condition. The joint loss gives $D$ two data constraints and makes
the Rho term identifiable from data rather than from assumption.

\subsection{Deferred to follow-up changes}

\begin{itemize}
\item \textbf{Joint $(\SWT, \SBCM)$ re-fit of $D$}. Motivated by the
  \texttt{yihO} failure and the endpoint-pileup cases.
\item \textbf{$\Delta$-metric at $n_{\text{runs}} \ge 5000$} or with
  paired-run variance reduction.
\item \textbf{Grid extension beyond $\Krut=0.13$}. Would catch genes
  where $\DstarWT$ \emph{under}-used Rho (argmin above 0.13 in the
  pointwise metric). Not observed in this panel but expected in the
  4273-gene genome-wide roll-out.
\item \textbf{Partial-BCM efficacy calibration}. If the genome-wide
  $K^{*,\text{pw}}_{\text{rut}}$ distribution sharp-peaks at some small
  positive value, that value is a single-number estimate of residual
  Rho activity under BCM in vivo.
\end{itemize}

\subsection{Known kernel artifact: ribosome snap-to-negative}
\label{sec:snap-artifact}

The ribosome elongation block in both \texttt{netseq\_tasep\_fast.py}
and the original \texttt{sjkimlab\_NETSEQ\_TASEP.m} contains a snap
rule that, on collision detection, sets
$\texttt{Ribo\_locs(RNA)} = \texttt{rnap\_locs(RNA)} - \texttt{RNAP\_width}$
(line 561 in MATLAB; line 595 in Python). When a ribosome loads at
position 1 with the RNAP at position 30 (the minimum required by the
\texttt{rnap\_locs >= 30} loading gate) and advances even one base on
the next time step, the overlap test fires and snaps the ribosome to
position $30 - 35 = -5$. The subsequent ribosome elongation block is
gated by \texttt{if Ribo\_locs > 0}, so the ribosome remains
\emph{permanently} stuck at the negative position --- it provides
no shielding for the rest of the RNAP's transit.

The Rho-loading test reads \texttt{Ribo\_locs} directly:
\begin{equation*}
  \texttt{PTRNAsize} \;=\; \texttt{rnap\_locs} - 65 - \texttt{Ribo\_locs},
\end{equation*}
so $\texttt{Ribo\_locs} = -5$ inflates the exposed-RNA length by 5\,bp
and crosses the $\texttt{PTRNAsize} > 50$ threshold five bases earlier
than the gap-based geometry would predict. Empirically the affected
RNAPs behave as if the ribosome were absent.

The snap fires whenever $\texttt{rnap\_locs} < 35 + \texttt{tempRibo}$
on the first elongation step after loading, which (for typical
$\texttt{tempRibo} \approx 2$) covers the regime
$\texttt{rnap\_locs}_{\text{at-load}} \lesssim 37$. The fraction of
RNAPs whose ribosome loads in this regime is approximately
$1 - \exp(-35\, k_{\text{ribo}} / v_{\text{RNAP}})$. For
\texttt{talB} ($k_{\text{ribo}} = 0.51$) this is $\sim$60\%; for
\texttt{insQ} ($k_{\text{ribo}} = 0.022$) it is $\sim$4\%. For
high-translation-efficiency genes a majority of RNAPs effectively
have no ribosome shielding --- this is the dominant explanation for
why kernel survival saturates well below the per-RNAP analytic
prediction (see \texttt{kinetic\_estimates\_Linfty\_xc.md} \S 9).

\paragraph{Implication for $D^*$.} The CMA-ES optimizes $D$ to make
the (artifacted) kernel reproduce $\SWT$. Whatever bias the artifact
introduces is absorbed into $D^*$ rather than into the rate
parameters (since $k_{\text{ribo}}$ is fixed at $\text{TE}/5$ and is
not co-fit). For \emph{kernel surrogate} use --- re-running flux
profiles, BCM sweeps, downstream analyses that go through the same
kernel --- $D^*$ is correct by construction. For \emph{kinetic
interpretation} as the absolute biological dwell time per bp, $D^*$
on high-$k_{\text{ribo}}$ genes is likely biased toward over-tall
pause peaks (the kernel needs extra decay-causing structure to match
the data when $\sim 60\%$ of its ribosomes are ineffective). Relative
peak structure --- where pauses sit, ratios of nearby peaks --- is
robust because the bias is approximately a global multiplicative
scaling on absolute amplitudes.

\paragraph{Status.} The artifact was in the canonical sjkimlab MATLAB
implementation, so all downstream scientific output that used the
unpatched kernel inherited it. The Python kernel (both
\texttt{netseq\_tasep\_fast.py} and \texttt{netseq\_tasep\_gpu.py})
has been patched with the clamp
$\texttt{Ribo\_locs} \leftarrow \max(1, \texttt{rnap\_locs} - \texttt{rnap\_width})$
at both snap sites (State~1 coupling-maintenance and State~4
free-ribosome collision). The 5-gene panel (\texttt{insQ},
\texttt{talB}, \texttt{aceA}, \texttt{dnaK}, \texttt{rpoB}) was
re-run with the patched kernel and the original (pre-patch) CMA-ES
$D^*$, with outputs in \texttt{ecoli/flux\_patched/}.

The diagnostic notebook \texttt{ribo\_bias.ipynb} reports the
per-gene bias the artifact pushed into $D^*$:

\begin{center}
\begin{tabular}{l r r r r}
\hline
gene & $k_{\text{ribo}}$ & $P_{\text{snap}}$ &
$\text{flux}_{\text{patched}}[L]$ & $\sqrt{\Delta_{\log\text{-MSE}}}$ \\
\hline
\texttt{insQ} & 0.022 & 0.04 & 0.183 & 0.05 \\
\texttt{rpoB} & 0.166 & 0.26 & 0.857 & 0.40 \\
\texttt{aceA} & 0.358 & 0.48 & 0.958 & 0.23 \\
\texttt{dnaK} & 0.456 & 0.57 & 0.991 & 0.54 \\
\texttt{talB} & 0.506 & 0.61 & 0.997 & 0.34 \\
\hline
\end{tabular}
\end{center}

where $\sqrt{\Delta} = \sqrt{\text{mean}((\log
\text{flux}_{\text{patched}} - \log
\text{flux}_{\text{orig}})^2)}$ is the per-gene log-RMS shift in
the patched-kernel survival relative to the original-kernel survival
that matches $S_{\text{exp}}$ by CMA-ES design.

\textbf{Implication for the existing $D^*$.} For high-$k_{\text{ribo}}$
genes (\texttt{talB}, \texttt{dnaK}), the patched-kernel survival at
the gene endpoint is $\approx$0.99 vs the original $\approx$0.4--0.5
that matches $S_{\text{exp}}$. The CMA-ES therefore inflated $D^*$
pause peaks by enough amplitude to recover the missing
$\approx$0.5 in survival via extra Rho-loading hazard. This bias is
$\sim 30$--$50\%$ in log-space ($\sqrt{\Delta} \approx 0.3$--$0.5$).
For low-$k_{\text{ribo}}$ \texttt{insQ}, the bias is negligible
($\sqrt{\Delta} \approx 0.05$, i.e.\ $\sim$5\% log-space shift) because
only $\approx$4\% of RNAPs were hit by the artifact.

The closed-form analytic M13 of \texttt{kinetic\_estimates\_Linfty\_xc.md}
\S 9.2 matches the patched-kernel flux endpoint within 5\% on 3 of 5
panel genes (\texttt{talB}, \texttt{aceA}, \texttt{dnaK}), confirming
M13 as the correct closed form for the fixed kernel.

Existing flux pickles in \texttt{ecoli/flux/} and the genome-wide
BCM sweep outputs were generated with the unpatched kernel and remain
internally self-consistent for kernel-surrogate use. A re-fit of CMA-ES
on the patched kernel (proposed as \texttt{cmaes-patched-rerun} if
warranted) would correct the absolute amplitude of $D^*$ pause peaks
on high-$k_{\text{ribo}}$ genes; relative pause structure (peak
positions, ratios within a gene) is robust to the artifact and
unchanged. See OpenSpec change \texttt{ribo-snap-bias-diagnostic}
and the \texttt{ribo\_bias.ipynb} notebook for details.

\section{CMA-ES deconvolution}
\label{sec:cmaes}

\subsection{Problem formulation}

For each gene of length $L$, we seek the intrinsic per-position dwell-time
profile $D^*(x) \in \mathbb{R}_{>0}^L$ such that the TASEP forward model,
when run with $D^*$ and the gene-specific kinetic constants, reproduces the
observed wildtype NET-seq density $S_{\text{exp,norm}}(x)$. We define the
objective as the mean-squared error between the simulated and observed
mean-normalized NET-seq profiles:
\begin{equation}
  \mathcal{L}(D) \;=\; \frac{1}{L}\sum_{x=1}^{L}
    \Bigl(\hat{S}_{\text{norm}}(D, x) - S_{\text{exp,norm}}(x)\Bigr)^2,
\end{equation}
where $\hat{S}_{\text{norm}}$ is the Monte Carlo estimate of the expected
NET-seq density under $D$, both sides divided by their respective spatial
means (so that the loss is scale-invariant and depends only on profile shape).

\subsection{Parameterization and initialization}

The search is conducted in log-space. Define
\begin{equation}
  \theta \;=\; \log D \;\in\; \mathbb{R}^L,
  \qquad D \;=\; \exp(\theta),
\end{equation}
where the exponential enforces positivity without constraints. Before each
kernel call, $D$ is mean-normalized,
\begin{equation}
  \tilde{D}(x) \;=\; \frac{D(x)}{\frac{1}{L}\sum_{x'} D(x')},
\end{equation}
so the mean dwell time is always 1 (in units of the baseline speed
$v_{\text{RNAP}} = 19\,\text{bp/s}$) and only the \emph{shape} of $D$ is
fitted.

The initial mean vector of the CMA-ES distribution is set directly from the
observed signal:
\begin{equation}
  \theta_0 \;=\; \log\bigl(\max(S_{\text{exp,norm}},\, 10^{-10})\bigr),
\end{equation}
i.e.\ the initial dwell profile is the experimental NET-seq profile itself.
This warm start places the optimizer close to a plausible solution from
generation zero, since $S_{\text{exp,norm}} \propto D \cdot j$ and $j$ is
close to 1 at low traffic density.

\subsection{CMA-ES with diagonal covariance}

We use the CMA-ES algorithm~\cite{hansen2016cma} as implemented in the
\texttt{pycma} library. For gene lengths $L \sim 300\text{--}4000$, the
standard full-covariance CMA-ES would require $O(L^2)$ storage and $O(L^3)$
covariance matrix updates per generation, which is prohibitive. We enable the
\texttt{CMA\_diagonal} option, which restricts the covariance matrix
$\mathbf{C}$ to its diagonal:
\begin{equation}
  \mathbf{C} \;=\; \operatorname{diag}(c_1^2,\, c_2^2,\, \ldots,\, c_L^2),
\end{equation}
reducing storage to $O(L)$ and per-generation update cost to $O(\lambda L)$
where $\lambda$ is the population size. This is well-suited to the TASEP
deconvolution problem because the per-position dwell times are largely
independent — Rho-mediated effects couple positions non-locally, but the
dominant signal in $D$ is the per-bp pause intensity, which is intrinsically
local.

The population size is set by the standard CMA-ES heuristic:
\begin{equation}
  \lambda \;=\; 4 + \lfloor 3 \ln L \rfloor,
\end{equation}
giving $\lambda = 21$ for a 300-bp gene and $\lambda = 25$ for a 1149-bp gene.
The initial step size is $\sigma_0 = 0.1$, corresponding to $\sim 10\%$
multiplicative perturbations in linear dwell space ($e^{\pm 0.1} \approx 1
\pm 0.1$). A fixed CMA seed (42) is used for reproducibility.

Each generation proceeds as:
\begin{enumerate}
  \item \textbf{Sample.} Draw $\lambda$ candidate vectors
    $\theta^{(1)}, \ldots, \theta^{(\lambda)} \sim \mathcal{N}(\mathbf{m},
    \sigma^2 \mathbf{C})$.

  \item \textbf{Evaluate.} For each candidate $\theta^{(i)}$, compute
    $D^{(i)} = \exp(\theta^{(i)})$, mean-normalize, and dispatch to the GPU
    kernel. The kernel runs $n_{\text{runs}}$ independent TASEP simulations
    in parallel (see \S\ref{sec:gpu}), averages the resulting NET-seq snapshot
    counts, mean-normalizes, and returns $\hat{S}_{\text{norm}}^{(i)}$. The
    MSE kernel computes $\mathcal{L}(\theta^{(i)})$ on-device.

  \item \textbf{Tell.} Pass the $\lambda$ fitness values back to the CMA-ES
    strategy (\texttt{es.tell}), which updates the mean $\mathbf{m}$, step
    size $\sigma$, and diagonal covariance $\mathbf{C}$ using the
    weighted recombination, cumulative step-size adaptation, and covariance
    matrix adaptation rules.
\end{enumerate}

\subsubsection*{Mean vector update}

The mean $\mathbf{m} \in \mathbb{R}^L$ is the current best estimate of
$\theta^*$ and the center of the sampling distribution. After evaluating all
$\lambda$ candidates, they are ranked by fitness (ascending MSE), the top
$\mu = \lfloor\lambda/2\rfloor$ are selected, and $\mathbf{m}$ is updated as
their weighted centroid:
\begin{equation}
  \mathbf{m}_{t+1} \;=\; \sum_{i=1}^{\mu} w_i\, \theta^{(i:\lambda)}_t,
\end{equation}
where $\theta^{(i:\lambda)}$ denotes the $i$-th ranked candidate (rank 1 =
lowest MSE) and the weights are the standard CMA-ES recombination weights:
\begin{equation}
  w_i \;=\; \frac{\ln(\mu + \tfrac{1}{2}) - \ln i}
                 {\displaystyle\sum_{j=1}^{\mu}\bigl[\ln(\mu+\tfrac{1}{2}) - \ln j\bigr]},
  \quad w_1 > w_2 > \cdots > w_\mu > 0, \quad \sum_{i=1}^{\mu} w_i = 1.
\end{equation}
For $\lambda = 25$ (insQ), $\mu = 12$ candidates are selected. The best
candidate receives $\approx 18\%$ of the total weight; the 12th receives
$\approx 1\%$. This soft-max selection avoids the brittleness of winner-takes-all
truncation while still biasing $\mathbf{m}$ strongly toward the highest-fitness
region.

The update is a pure weighted mean — $\mathbf{m}$ never evaluates a gradient
directly. It infers the direction of improvement solely from the rank order of
fitness values across the sampled cloud, which is why CMA-ES is robust to the
stochastic noise introduced by finite $n_{\text{runs}}$ Monte Carlo estimation.

The warm start $\mathbf{m}_0 = \log(S_{\text{exp,norm}})$ is consequential:
since $\mathbf{m}$ drifts by weighted recombination, the optimizer begins
sampling around a plausible initial dwell profile. If $\mathbf{m}_0$ is already
close to $\theta^*$, the initial cloud of candidates contains solutions near
the optimum and $\mathbf{m}$ converges faster than it would from a random
initialization.

\subsubsection*{Why log-space makes the update well-scaled}

In log-space, a perturbation $\delta\theta$ corresponds to a multiplicative
change $e^{\delta\theta} \approx 1 + \delta\theta$ in linear dwell. The
per-position coordinate scales are commensurate regardless of the absolute
magnitude of $D^*(x)$: a fast site with $D^*(x) \approx 0.01$ and a strong
pause site with $D^*(x) \approx 10$ both respond to a step of $\sigma = 0.1$
as a $\sim 10\%$ multiplicative perturbation. In linear space, the same fixed
step size would be a $10\times$ change at the fast site and a $1\%$ change at
the pause — the weighted mean $\mathbf{m}$ would be dominated by the
high-magnitude pause coordinates at every update step.

\subsection{Multi-fidelity $n_{\text{runs}}$ schedule}

Evaluating $\mathcal{L}$ via Monte Carlo introduces fitness noise proportional
to $1/\sqrt{n_{\text{runs}}}$. More runs reduce noise but increase the GPU
kernel launch time linearly up to the SM-pass boundary (see Table~\ref{tab:occupancy}).
In early generations (large $\sigma$, coarse landscape), high fitness noise is
tolerable because only the ranking of candidates matters, not precise MSE values.
As $\sigma$ contracts and the population converges, finer MSE discrimination is
needed.

The \texttt{NRunsSchedule} implements this as a sigma-ratio schedule:
\begin{equation}
  n_{\text{runs}}(\sigma) \;=\;
  \begin{cases}
    n_{\text{low}}  & \text{if } \sigma/\sigma_0 > 0.8, \\
    n_{\text{med}}  & \text{if } 0.6 < \sigma/\sigma_0 \le 0.8, \\
    n_{\text{high}} & \text{if } \sigma/\sigma_0 \le 0.6.
  \end{cases}
\end{equation}
The operational setting for the genome-wide run uses the uniform schedule
$n_{\text{low}} = n_{\text{med}} = n_{\text{high}} = 200$, keeping all kernel
launches within a single SM pass ($\sim 2.3\,\text{s/gen}$). This trades
$\sim 5\%$ final MSE for a $19\%$ wall-time reduction relative to the
escalating $200/400/800$ schedule.

\subsection{Early stopping}

To avoid spending generations on converged fits, a sliding-window criterion
halts optimization when the relative improvement in the best-seen fitness falls
below a threshold over a fixed window:
\begin{equation}
  \frac{\mathcal{L}_{\text{best}}(g - W) - \mathcal{L}_{\text{best}}(g)}
       {\mathcal{L}_{\text{best}}(g - W)}
  \;\;<\;\; \epsilon_{\text{stop}},
\end{equation}
with $W = 30$ generations and $\epsilon_{\text{stop}} = 0.005$ (0.5\%).
On the 5-gene benchmark, 2 of 5 genes triggered early stopping (dnaK at
gen~127, rpoB at gen~184 of a 200-gen budget), reducing their wall time
without material loss of fit quality.

\subsection{Output and checkpointing}

Every 25 generations the optimizer state is checkpointed to
\texttt{cmaes\_checkpoints/gen\_\{N\}.pkl}, enabling resume after interruption.
On convergence the best-ever solution $\theta^*$ is extracted (\texttt{es.result.xbest}),
exponentiated, and mean-normalized:
\begin{equation}
  D^*(x) \;=\; \frac{\exp(\theta^*(x))}{\frac{1}{L}\sum_{x'}\exp(\theta^*(x'))}.
\end{equation}
The per-gene output pickle (\texttt{ecoli/cmaes/\{gene\}.pkl}) stores
$D^*$, $\theta^*$, the reference signal $S_{\text{exp,norm}}$, the per-gene
kinetic constants ($\kribo$, $\Krut$), and the optimization history
(best MSE per generation, $\sigma$ trajectory, wall time).

\section{RNAP flux from $D^*$}
\label{sec:flux}

\subsection{Definition}

The RNAP flux $j(x)$ at position $x$ is the mean number of RNAPs that
traverse position $x$ per unit simulation time at steady state. It is
the throughput analogue of the local density: whereas
$\rho(x) = D(x)\,j(x)$ counts how many RNAPs \emph{dwell at} $x$, $j(x)$
counts how many \emph{pass through} $x$. The two quantities together
decompose the NET-seq signal into an intrinsic pause component ($D^*$)
and a throughput component ($j$).

\subsection{Simulation-based estimation}

Given $D^*(x)$ from the CMA-ES deconvolution, flux is estimated by
running $n_{\text{runs}}$ independent TASEP simulations and counting
RNAP traversals. In the GPU kernel, every time an RNAP advances from
position $x_{\text{curr}}$ by $\Delta x_{\text{adv}}$ bases, a counter
is incremented atomically for each position traversed:
\begin{equation}
  \texttt{flux\_count}[x] \;\mathrel{+}=\; 1
  \quad \forall\, x \in \bigl\{x_{\text{curr}},\,\ldots,\,
                               x_{\text{curr}} + \Delta x_{\text{adv}} - 1\bigr\}.
\end{equation}
After all $n_{\text{runs}}$ simulations the raw count is averaged and
normalized to the promoter-proximal position ($x=1$):
\begin{equation}
  \hat{j}(x) \;=\; \frac{\texttt{flux\_count}[x]}{n_{\text{runs}}},
  \qquad
  j_{\text{norm}}(x) \;=\; \frac{\hat{j}(x)}{\hat{j}(1)}.
\end{equation}
$j_{\text{norm}}(x)$ is the fraction of promoter-initiated RNAPs that
successfully reach position $x$. A value $j_{\text{norm}}(x) < 1$
reflects cumulative attrition from Rho-mediated premature termination and
hard-body collision-induced stalling upstream of $x$.

\subsection{Relationship to $D^*$ and the NET-seq signal}

At steady state the NET-seq density factorizes as (see \S\ref{sec:math}):
\begin{equation}
  S(x) \;\propto\; D^*(x)\cdot j(x;\,D^*,\,\theta).
\end{equation}
The CMA-ES fit minimized the MSE on $S(x)$; flux is a derived quantity
not directly optimized. Computing $j_{\text{norm}}(x)$ post-hoc from
$D^*$ therefore provides an independent view that is not available from
$D^*$ alone: it directly quantifies how much Rho termination attenuates
throughput between the TSS and any downstream position, and separates
intrinsic pausing from transcript survival.

\subsection{Practical computation}

Flux is computed by \texttt{postprocess\_flux.py}, which reads each
gene's CMA-ES pickle (\texttt{ecoli/cmaes/\{gene\}.pkl}), reconstructs
$D^*$ from $\theta^*$, and calls \texttt{simulate\_with\_flux\_gpu} with
\texttt{compute\_flux=1} (a flag that activates the traversal counters;
it is disabled during CMA-ES optimization to avoid the per-step atomic
overhead). Runs use $dt = 0.1$ and $n_{\text{runs}} = 200$.

The per-gene output (\texttt{ecoli/flux/\{gene\}\_flux.pkl}) stores
$j_{\text{norm}}$ as a $(L,)$ float64 array. Together with the
CMA-ES pickle's $D^*$, this provides the full steady-state
characterization of the gene under the fitted model:
\begin{equation}
  S_{\text{sim}}(x) \;\propto\; D^*(x)\cdot j_{\text{norm}}(x),
\end{equation}
where the proportionality constant absorbs the absolute loading rate
$\kload$ and simulation duration.

\section{Analytic approximation for the flux}
\label{sec:flux_analytic}

\subsection{Continuity-equation Gaussian}

The flux $j(x)$ from \S\ref{sec:flux} obeys the mean-field continuity
equation
\begin{equation}
  \frac{\mathrm{d}j}{\mathrm{d}x} \;=\; -\gamma(x)\, j(x),
  \qquad
  j_{\text{norm}}(x) \;=\; \exp\!\left(-\int_0^x \gamma(x')\,\mathrm{d}x'\right),
\end{equation}
where $\gamma(x)$ is the per-bp Rho-termination rate. The kernel performs
one Rho-loading Bernoulli draw per timestep with success probability
$K_{\text{rut}}\, L_{\text{unc}}(x) / L_{\text{gene}}$, while the RNAP
advances on the order of $v_{\text{RNAP}} \cdot \mathrm{d}t$ bp per timestep.
Converting per-timestep to per-bp gives
\begin{equation}
  \gamma(x) \;=\; \frac{K_{\text{rut}}\, L_{\text{unc}}(x)}
                       {v_{\text{RNAP}}\, L_{\text{gene}}},
  \qquad v_{\text{RNAP}} = 19\ \text{bp/s}.
  \label{eq:gamma}
\end{equation}
In the \emph{strict-uncoupled} limit ($x_{\text{ribo}} \approx 0$) the
uncovered length grows linearly, $L_{\text{unc}}(x) = x - 65$ for
$x > \ell_{\text{load}}$, with $65 = w_{\text{RNAP}} + 30$ and
$\ell_{\text{load}} = 65 + 50 = 115$ (the kernel requires
$L_{\text{unc}} > 50$ to trigger Rho loading). Integrating
\eqref{eq:gamma} gives a Gaussian survival profile:
\begin{equation}
  \Pi_\rho^{\text{uncoupled}}(x) \;=\;
    \exp\!\left(-\frac{K_{\text{rut}}}{2\, v_{\text{RNAP}}\, L_{\text{gene}}}
                       \,(x - 65)^2\right),
  \qquad x > \ell_{\text{load}}.
  \label{eq:gauss-uncoupled}
\end{equation}

\subsection{Empirical shielding factor}

\paragraph{Normalization.} We normalize $j(x)$ to its value at $x=1$,
$j_{n0}(x) := j(x)/j(1)$, so that $j_{n0}$ is directly the
\emph{survival fraction}: the probability that a promoter-initiated
RNAP has not yet terminated by position $x$. Under this normalization,
\eqref{eq:gauss-uncoupled} matches $j_{n0}$ position-for-position with
no rescaling at the boundary, and the residual after dividing out the
Gaussian fit is pinned at $1$ at $x=1$ by construction.

The strict-uncoupled limit overestimates Rho attrition for all panel
genes: the simulated $j_{n0}(L)$ is in $0.16$--$0.65$, whereas
\eqref{eq:gauss-uncoupled} would predict $10^{-6}$ to $0.03$ depending
on $L$. A single empirical shielding factor $s \in (0,1]$ multiplied
into the Gaussian exponent matches $j_{n0}(x)$ on every panel gene to
within $\sim 5$--13\% RMS in log-space:
\begin{equation}
  \Pi_\rho^{\text{fit}}(x) \;=\;
    \exp\!\left(-s \cdot \frac{K_{\text{rut}}}
                                  {2\, v_{\text{RNAP}}\, L_{\text{gene}}}
                       \,(x - 65)^2\right).
  \label{eq:gauss-fit}
\end{equation}
$s$ has a direct physical reading: $s \cdot (x - 65)$ is the
\emph{effective uncovered-RNA length} along the gene, set by the
steady-state ribosome lag behind the RNAP. $s = 1$ corresponds to a
fully uncoupled gene; $s \to 0$ means the ribosome tracks the RNAP at
the minimum coupling distance.

\begin{table}[h]
\centering
\small
\begin{tabular}{lrrrrr}
\toprule
Gene & $L$ (bp) & $\kribo$ & Shield $s$ & $j_{n0}(L)$ & RMS log \\
\midrule
insQ & 1149 & 0.022 & 0.53 & 0.158 & 0.087 \\
talB &  954 & 0.506 & 0.25 & 0.518 & 0.039 \\
aceA & 1305 & 0.358 & 0.12 & 0.653 & 0.045 \\
dnaK & 1917 & 0.456 & 0.16 & 0.433 & 0.127 \\
rpoB & 4029 & 0.166 & 0.05 & 0.500 & 0.123 \\
\bottomrule
\end{tabular}
\caption{Single-parameter Gaussian fit \eqref{eq:gauss-fit} to
$j_{n0}(x) = j(x)/j(1)$ from \texttt{ecoli/flux/\{gene\}\_flux.pkl}.
$j_{n0}(L)$ is the survival fraction at the 3$^\prime$ end. RMS log is
the root-mean-square log-space deviation between simulated and fitted
survival. All 5 panel genes match within $<14\%$ log-deviation by a
one-parameter Gaussian, with no traffic correction.}
\label{tab:shield}
\end{table}

\paragraph{Anticorrelation with translation efficiency.} The smallest
shield ($s = 0.05$) belongs to \texttt{rpoB} and the largest ($s = 0.54$)
to \texttt{insQ}, which has the lowest ribosome loading rate
($\kribo = 0.022$). The intuition is consistent: when ribosomes are
scarce, the RNAP runs ahead and exposed RNA approaches the uncoupled
$(x-65)$ limit; when ribosomes load fast, they track the RNAP and clamp
$L_{\text{unc}}$ near its minimum threshold. A predictive closed-form
for $s$ as a function of $\kribo$ (and $v_{\text{ribo}}$,
$v_{\text{RNAP}}$) is a natural next step.

\subsection{The traffic throughput term}

The remaining factor in the full flux decomposition (see
\S\ref{sec:math}),
\begin{equation}
  j(x) \;=\; \kload \cdot \Pi_\rho(x) \cdot \Pi_{\text{traffic}}(x),
\end{equation}
is the throughput modifier from hard-body exclusion. The residual
$j_{n0}(x) / \Pi_\rho^{\text{fit}}(x)$ acts as a proxy for
$\Pi_{\text{traffic}}(x)$, pinned at $1$ at $x=1$ by construction and
empirically within $\pm 20\%$ of unity across the gene. The residual
splits into two qualitatively different components:

\paragraph{Broad positive bias.} Through the middle of each gene the
residual sits slightly above $1$ (e.g., $1.05$--$1.20$ for \texttt{insQ}
in $x \in [600,\,850]$). This is a \emph{$\Pi_\rho$ shape mismatch},
not traffic: the true $\gamma(x)$ plateaus once the ribosome reaches
its steady-state lag behind the RNAP, whereas the Gaussian assumes
$L_{\text{unc}}(x) = s\cdot(x-65)$ grows without bound. A refined
ansatz
\begin{equation}
  \gamma(x) \;=\; \gamma_\infty\bigl(1 - e^{-(x-\ell_{\text{load}})/x_*}\bigr)
\end{equation}
with equilibration length $x_*$ and saturated rate $\gamma_\infty$
would absorb this systematic bias.

\paragraph{Localized dips.} Sharp negative excursions in the residual
coincide with strong $D^*$ pause sites (e.g., \texttt{insQ} drops to
$0.78$ at $x \approx 940$, near a high-dwell feature). These are the
genuine traffic signature. The expected analytic form in the
low-density TASEP limit is
\begin{equation}
  \Pi_{\text{traffic}}(x) \;\approx\; 1 \;-\; a \cdot \rho(x+1),
  \qquad \rho(x) \;=\; D^*(x)\, j(x),
  \label{eq:traffic}
\end{equation}
with $a \sim w_{\text{RNAP}} / L_{\text{gene}}$ a per-site exclusion
volume. Validating the residual dips against the dominant peaks of
$D^*$ and fitting $a$ would close the loop on the full analytic flux
model.

\paragraph{Practical use.} The Gaussian fit gives an immediate
sanity check on $j_{\text{norm}}$ output from the kernel: anomalously
large RMS log deviation flags genes where the scalar-$K_{\text{rut}}$
model has structural problems (e.g., the \texttt{yihO}-class
model failures from \S\ref{sec:results}). The analytic form
$\Pi_\rho^{\text{fit}}$ also offers a fast proxy for genome-wide
analyses that would otherwise require running the simulator end-to-end.
Implementation and 5-gene validation are in
\texttt{mathmodel-netseq.ipynb}.

\section{Computational acceleration architecture}
\label{sec:gpu}

\subsection{The bottleneck: sequential stochastic simulation at genome scale}

The TASEP forward model is a discrete-time stochastic simulation that must be
run sequentially in time: at each timestep $\Delta t = 0.1\,\text{s}$, every
RNAP, ribosome, and Rho complex on the gene is advanced according to
first-passage dwell-time draws, subject to hard-body exclusion constraints.
A single simulation run for a gene of length $L$ takes
$T_{\text{sim}}/\Delta t \approx 20{,}000$ timesteps; at each step the inner
loop touches $O(n_{\text{RNAP}}) \approx 30$ particles.

CMA-ES evaluation requires $n_{\text{pop}} \approx 25$ candidates, each
evaluated at $n_{\text{runs}} = 200$ independent simulation runs to estimate
the expected NET-seq signal. A single CMA-ES generation for a $\sim$1000-bp
gene therefore performs $25 \times 200 \times 20{,}000 \times 30 \approx 3
\times 10^9$ stochastic particle updates, each gated by a random
exponential draw and conditional collision check. Across 4273 E.\ coli genes
at $\sim$200 generations each, the genome-wide problem entails $\sim
10^{16}$ stochastic operations total.

The simulation is branch-heavy and latency-dominated: the inner loop
branches on particle state, collision checks, Rho-loading decisions, and
ribosome-coupling state. SIMD throughput and FP32/FP64 pipelines are
irrelevant bottlenecks. The accessible speedup strategy is massive thread
parallelism across independent simulation runs, not arithmetic acceleration.

\subsection{GPU threading model}

The GPU kernel (\texttt{\_tasep\_core\_cuda}, \texttt{netseq\_tasep\_gpu.py})
maps one CUDA thread to one simulation run. A launch of
$N_{\text{total}} = n_{\text{pop}} \times n_{\text{runs}}$ threads
(e.g.\ $25 \times 200 = 5000$) evaluates one full CMA-ES generation in a
single kernel call. Each thread runs the complete sequential TASEP loop
independently, with no inter-thread communication during the simulation.

The thread index encodes both the candidate and the run:
\begin{equation}
  c_{\text{idx}} = \lfloor t_{\text{id}} / n_{\text{runs}} \rfloor,
  \qquad
  r_{\text{idx}} = t_{\text{id}} \bmod n_{\text{runs}}.
\end{equation}
Thread $t_{\text{id}}$ reads the dwell profile of candidate $c_{\text{idx}}$,
draws its own independent RNG stream (xoroshiro128p), and writes its NET-seq
snapshot counts atomically to a shared output buffer indexed by
$(c_{\text{idx}}, x)$.

\subsubsection*{Thread-local state}

Each thread maintains five HIST\_CAP $= 128$ element arrays in
\emph{local (register-spill) memory}:

\begin{center}
\small
\begin{tabular}{lll}
\toprule
Array & dtype & Purpose \\
\midrule
\texttt{rnap\_locs[128]} & int32 & per-RNAP position \\
\texttt{ribo\_locs[128]} & int32 & per-RNAP ribosome position \\
\texttt{rho\_locs[128]}  & int32 & per-RNAP Rho position \\
\texttt{rnap\_ribo\_coupling[128]} & int32 & coupling state flag \\
\texttt{ribo\_loadt[128]}  & float32 & ribosome load time \\
\texttt{exit\_buf[32]}    & float32 & first-passage scratch \\
\bottomrule
\end{tabular}
\end{center}

Total local memory per thread: $\approx 2.7\,\text{KB}$ (v3.0, float32/int32).
The prior version (v2.2) used float64/int64, giving $\approx 5.3\,\text{KB}$
per thread. The v3.0 reduction to float32 and int32 halves local memory,
slightly improving L1 cache locality, though it does not change per-generation
wall time because the bottleneck is instruction latency, not memory bandwidth.

\subsubsection*{Simulation loop structure}

The main loop iterates over $n_{\text{steps}} = T_{\text{sim}}/\Delta t$
timesteps and, at each step, executes four sub-steps in order, mirroring the
CPU reference exactly:

\begin{enumerate}
  \item \textbf{RNAP loading.} Draw next load time from
    $\text{Exp}(1/k_{\text{load}})$. Block if last loaded RNAP has not
    cleared its footprint ($w_{\text{RNAP}} = 35\,\text{bp}$).

  \item \textbf{RNAP elongation.} For each active RNAP, call
    \texttt{\_sample\_first\_passage\_cuda} to draw the number of bases
    advanced in this timestep from the per-position dwell profile; apply
    hard-body collision check against the next-upstream RNAP.

  \item \textbf{Rho loading and elongation.} Per active RNAP, compute the
    uncovered RNA length $L_{\text{unc}} = x_{\text{RNAP}} - w_{\text{RNAP}}
    - 30 - x_{\text{ribo}}$; if $L_{\text{unc}} > 50\,\text{nt}$, draw a
    Bernoulli with probability $K_{\text{rut}} \cdot L_{\text{unc}} /
    L_{\text{gene}} \cdot 100\,\Delta t$ to load Rho at a uniform random
    position on the exposed RNA. Advance Rho along the rho-dwell profile
    (speed $= 5 \times v_{\text{ribo}}$). If Rho catches the RNAP, mark
    the RNAP as prematurely terminated.

  \item \textbf{Ribosome elongation.} Advance each ribosome from its
    coupled or free state; detect RNAP--ribosome coupling on collision.
\end{enumerate}

At seven fixed snapshot times (200, 400, 600, 800, 1000, 1200, 1400 s),
each thread records its RNAP positions atomically into the
$(n_{\text{pop}}, L)$ output buffer:
\begin{equation}
  \texttt{out\_netseq\_sum}[c_{\text{idx}},\, x-2]
  \;\mathrel{+}=\; 1
  \quad \forall \text{ RNAP at position } x.
\end{equation}
A \texttt{first\_active} index is maintained to skip RNAPs that have
already terminated or exited the gene, matching the CPU reference
implementation exactly.

\subsection{Dwell profile deduplication}

In the naive (v2.2) implementation, the host replicated each candidate's dwell
profile $n_{\text{runs}}$ times before transfer, producing a
$(n_{\text{total}},\, L+1)$ device array of size $\sim 92\,\text{MB}$ for a
typical CMA-ES generation on insQ ($n_{\text{total}} = 10{,}000,\, L = 1149$).

v3.0 deduplicates: the kernel receives a $(n_{\text{pop}},\, L+1)$ array
and each thread reads \texttt{dwell\_profiles[candidate\_idx, :]}. The
$n_{\text{runs}}$ threads sharing the same candidate share the same dwell
row via the L2 cache. Memory reduction:
\begin{equation}
  92\,\text{MB} \;\longrightarrow\; 0.12\,\text{MB} \quad (\text{insQ},
  \; n_{\text{pop}}=25, \; L=1149, \; \text{float32}).
\end{equation}
The per-candidate ribosome loading rate \texttt{k\_ribo\_loading\_arr} and
Rho loading prefactor \texttt{pt\_percent\_arr} follow the same pattern:
$(n_{\text{pop}},)$ device arrays indexed by $c_{\text{idx}}$.

\subsection{Atomic output reduction}

v2.2 allocated a $(n_{\text{total}},\, L)$ output buffer ($\sim 46\,\text{MB}$
for insQ) so each thread could write its NET-seq snapshot without contention,
followed by a host-side summation over runs per candidate.

v3.0 threads instead write directly to a $(n_{\text{pop}},\, L)$ buffer
($\sim 0.11\,\text{MB}$) via \texttt{cuda.atomic.add}. On Ampere GPUs, L2
hardware atomics are efficient: with $n_{\text{runs}} = 200$ threads per
candidate writing to $L \approx 1000$ positions, the mean number of
concurrent writes per address per snapshot is $\sim 0.04$ — essentially
contention-free. Total output transfer to host is reduced 1380$\times$.

\subsection{On-device MSE kernel}

A second kernel, \texttt{\_compute\_mse\_cuda}, eliminates the host-side MSE
computation entirely. It is launched as $n_{\text{pop}}$ blocks of 256
threads, one block per candidate:

\begin{enumerate}
  \item \textbf{Pass 1 — normalization sum.} Each thread accumulates
    \texttt{netseq\_sum[c\_idx, pos]} over a stride-256 tile of positions
    into shared memory, followed by a parallel binary-tree reduction across
    the 256 threads to obtain the total count $S_{\text{total}}$.

  \item \textbf{Pass 2 — MSE.} Using $s_{\text{mean}} =
    S_{\text{total}} / (n_{\text{runs}} \cdot L)$ for normalization, each
    thread accumulates $\bigl(S_{\text{exp}}(x) -
    \hat{S}_{\text{sim}}(x)\bigr)^2$ over its tile; another binary-tree
    reduction yields $\sum_x (\cdot)^2 / L$, written to
    \texttt{mse\_out[c\_idx]}.
\end{enumerate}

The host receives only an $n_{\text{pop}}$-float array ($\sim$100 bytes)
per generation instead of $n_{\text{pop}} \times L$ floats ($\sim$0.11 MB).

\subsection{Performance profile and bottleneck analysis}

\begin{table}[h]
\centering
\small
\begin{tabular}{rrrr}
\toprule
$n_{\text{runs}}$ & Blocks (RTX A5000) & SM passes & Time/gen (insQ) \\
\midrule
100 & 20 & 1 & 2.26\,s \\
200 & 40 & 1 & 2.29\,s \\
300 & 59 & 1 & 2.28\,s \\
400 & 79 & 2 & 4.35\,s \\
800 & 157 & 3 & 6.9\,s \\
\bottomrule
\end{tabular}
\caption{Per-generation wall time on one RTX A5000 (64 SMs) as a function of
$n_{\text{runs}}$ ($n_{\text{pop}} = 25$, $L = 1149\,\text{bp}$, insQ).
Time is flat within each SM-pass band, then jumps discretely at each SM-pass
boundary.}
\label{tab:occupancy}
\end{table}

The kernel is \emph{latency-bound}: per-generation time is governed by
$\lceil N_{\text{blocks}} / N_{\text{SMs}} \rceil \times t_{\text{thread}}$,
where $t_{\text{thread}}$ scales linearly with gene length and is set by
instruction-issue latency, not arithmetic throughput. The RTX A5000 has 64
SMs; a launch of 40 blocks (n\_runs=200, n\_pop=25 with 128 threads/block)
fills all 64 SMs in a single pass ($\sim 2.3\,\text{s}$). Increasing
$n_{\text{runs}}$ from 200 to 400 doubles the block count to 79,
requiring a second SM pass and doubling the wall time for no reduction in
per-generation cost.

Several optimization strategies that would help compute-bound kernels were
tested and yielded no benefit here:

\begin{itemize}
  \item \textbf{Float32 conversion.} Float32 arithmetic is 32$\times$ faster
    than float64 in FP-throughput terms on Ampere. The TASEP kernel is branch-
    and latency-bound; removing float64 did not change per-generation time.
    Float32 was retained for the 2$\times$ local-memory reduction it provides.

  \item \textbf{Concurrent CUDA streams.} Launching two or three genes'
    kernels concurrently in separate streams was expected to hide latency via
    interleaving. Benchmarked result: 1.00$\times$ speedup (exact serialization).
    Root cause: each single-gene kernel already fills all 64 SMs; concurrent
    streams interleave blocks, making each gene take $N\times$ longer while
    throughput stays flat.

  \item \textbf{Common Random Numbers (CRN).} Pre-tiling the same $n_{\text{runs}}$
    RNG states across candidates was implemented and confirmed to give zero
    variance for identical $\theta$ vectors. However, CMA-ES ranking variance
    was not reduced (46.84 with CRN vs.\ 46.50 independent), because TASEP's
    state-dependent first-passage sampling consumes different numbers of random
    draws per step for different dwell profiles, desynchronizing the shared
    streams after the first few steps.

  \item \textbf{Surrogate-assisted CMA-ES.} Random Forest and Gaussian Process
    surrogates were trained on 750 CMA-ES samples (generations 0--29) to
    predict candidate fitness. All models produced negative $R^2$ on the
    test set (gens 30--49): the CMA-ES search trajectory is non-stationary, and
    the $1149$-dimensional parameter space is too sparse for a lightweight
    surrogate.
\end{itemize}

The actionable lever is keeping $n_{\text{runs}} \le 300$ to stay within one
SM pass. The operational choice of $n_{\text{runs}} = 200$ gives
$\sim 2.3\,\text{s/gen}$, the minimum achievable on this hardware at this
population size.

\subsection{GPU vs.\ CPU and hardware scaling}

Per-generation time scales approximately linearly with gene length:
\begin{equation}
  t_{\text{gen}}(L) \;\approx\; 1.524 \times 10^{-3} \cdot L + 0.535 \;\text{s}.
\end{equation}
For insQ ($L = 1149\,\text{bp}$) at $n_{\text{runs}} = 200$, the RTX A5000
achieves $2.29\,\text{s/gen}$ vs.\ $4.88\,\text{s/gen}$ on a 32-thread CPU
(Threadripper PRO): a 2.1$\times$ GPU speedup. This modest ratio reflects the
branch-heavy, latency-dominated character of the kernel: the GPU's advantage
comes entirely from its massive thread count, not from its arithmetic pipeline.

Because the kernel is clock-speed- and SM-count-limited (not compute-limited),
GPU generation is counterintuitive: an A100 with 108 SMs and a 1.41 GHz boost
clock is $\sim$20\% \emph{slower} than the RTX A5000 (64 SMs, 1.70 GHz) at
$n_{\text{runs}} = 200$, because both fit all blocks in one SM pass and the
A100's lower clock speed dominates. Only the H100 (1.83 GHz) or H200 (1.98
GHz) improve on the A5000, by 8--17\%.

\subsection{CMA-ES integration and genome-wide runner}

Each gene is deconvolved independently by \texttt{run\_cmaes\_for\_gene()}
(\texttt{cmaes\_multifidelity.py}). The optimizer is initialized from the
log-normalized experimental NET-seq signal:
\begin{equation}
  \theta_0 = \log\bigl(\max(S_{\text{exp,norm}},\, 10^{-10})\bigr),
\end{equation}
giving a warm start that places the initial dwell profile close to the
observed density shape. The per-gene ribosome loading rate is set from the
translation efficiency table as $k_{\text{ribo}} = \text{TE}/5$.

The CMA-ES population size scales with gene length ($n_{\text{pop}} \approx
4 + \lfloor 3 \ln L \rfloor$). Fitness noise from finite $n_{\text{runs}}$
is absorbed by the CMA-ES noise-handling machinery (noise rank update, noise
epsilon injection).

The genome-wide runner (\texttt{run\_genome\_wide.py}) assigns one OS process
per GPU via \texttt{multiprocessing.spawn} to avoid GIL contention: CMA-ES
operations (\texttt{es.ask()}, \texttt{es.tell()}) hold the GIL, and a
threading model would starve the second GPU's kernel launches. With independent
processes, two GPUs operate in full parallel, achieving a 2$\times$ throughput
on the 4273-gene corpus. Genes are sorted by descending length before dispatch
so that the longest (slowest) genes run first, minimizing tail-latency imbalance
at the end of the run.

\section{Reproducibility}

\begin{lstlisting}[basicstyle=\ttfamily\small, breaklines=true]
python bcm_sweep.py --results-dir ecoli \
  --genes aceA,dnaK,rpoB,talB,insQ,rhsA,rhsB,ydbA,yihO,caiT \
  --device-id 0 --n-runs 5000
\end{lstlisting}

Panel wall time: $1472\,\text{s}$ on a single RTX A5000. Per-gene
output: \texttt{ecoli/bcm\_sweep/\{gene\}\_sweep.pkl}. Analysis
notebook: \texttt{bcm\_sweep\_analysis.ipynb}. OpenSpec change
\texttt{openspec/changes/bcm-krut-sweep} captures the full design and
task list (23/31 complete; stages 5--6 are the genome-wide roll-out).

\end{document}
